% NinuxOS_Architecture.tex - the system and the files of the project, as of v2.9.60.
% Written anew on 10 October 2026 from the running system and its disk image; the old
% document was lost with the disk.   pdflatex NinuxOS_Architecture.tex  (twice, for the contents)
\documentclass[11pt,a4paper]{article}
\usepackage[T1]{fontenc}
\usepackage[utf8]{inputenc}
\usepackage[margin=2.4cm]{geometry}
\usepackage{booktabs}
\usepackage{longtable}
\usepackage{array}
\usepackage{enumitem}
\usepackage[hidelinks]{hyperref}

\newcommand{\file}[1]{\texttt{#1}}
\newcommand{\cmd}[1]{\texttt{#1}}
\setlist{itemsep=2pt,topsep=4pt}

\title{NinuxOS\\[4pt]\large Architecture of the system and of the project\\[2pt]
\normalsize version 2.9.60}
\author{}
\date{10 October 2026}

\begin{document}
\maketitle

\noindent\textbf{How to read this.} Everything here was read from the system itself: from the disk
image of v2.9.59, from the objects in its cabinet, from the text of the kernel and from the system
running in QEMU. Where something is not known yet, it says so; section~\ref{sec:open} lists those
places. Nothing is taken from memory of the old documents.

\tableofcontents
\newpage

% ------------------------------------------------------------------------------------------
\section{The idea}

NinuxOS is an operating system of \emph{functional objects}. Its own start line says it:
\begin{quote}\ttfamily NinuxOS v2.9.59 - Functional Objects (DATA + META + SOCKETS)\end{quote}

\begin{itemize}
\item \textbf{Everything is an object.} A program, a picture, a page, a theme, a text, a drawer,
      the machine's identity, the log: each is one ASMO object. An object carries its \emph{data},
      its \emph{meta data} (facts as \texttt{key:value}) and its \emph{sockets} (its own functions).
      A picture can carry the code that shows it; a program is an object whose sockets do the work.
      This is programming \emph{with} objects, not object-oriented programming: there are no classes.
\item \textbf{No file system.} Objects lie in the \emph{cabinet}, one after another on the disk;
      \emph{drawers} hold more objects. Formats such as FAT or NTFS are only read from foreign
      sticks (\cmd{files}).
\item \textbf{Two kinds of object files.} \file{.asmo} is an object of the system or of the user;
      \file{.kasm} is a part of the kernel (the kernel itself and the \emph{kernel suits} it puts
      on when needed).
\item \textbf{One language for programs.} Programs are written in PyTon, a dialect with the syntax
      of Python. It is translated to x86-64 assembly and stored as the code of an ASMO object. The
      compiler is inside the kernel: the system builds its own programs (\cmd{build}, Studio).
\item \textbf{Pages.} Interfaces that are not programs are HTML with CSS and PyTon scripts, shown
      by NET (the browser of the system): the start menu, the taskbar, the desktop menu, Setup.
\end{itemize}

% ------------------------------------------------------------------------------------------
\section{The ASMO v4 object}\label{sec:asmo}

An object is one run of bytes. All numbers are little-endian; 8 bytes wide unless said otherwise.

\subsection{The header (0x68 bytes)}

\begin{longtable}{@{}l p{11.5cm}@{}}
\toprule
offset & meaning \\ \midrule
\texttt{0x00} & the signature, 4 bytes: \texttt{4F 4D 53 41} (the text \texttt{OMSA}) \\
\texttt{0x04} & the version of the format, 4 bytes: 4 \\
\texttt{0x08} & the size of the whole object \\
\texttt{0x10} & where DATA starts (always \texttt{0x68}) \\
\texttt{0x18} & where META starts \\
\texttt{0x20} & where the table of the system sockets starts \\
\texttt{0x28} & where the table of the named sockets starts \\
\texttt{0x30} & where the table of the prockets starts \\
\texttt{0x38} & 1: the object has code; 0: it has none (a source, a drawer) \\
\texttt{0x40} & 0 in every object seen \\
\texttt{0x48} & 4 bytes: 1 or 2 (not understood; a generation or a state) \\
\texttt{0x4C} & 4 bytes: the identity of the object; \cmd{id NAME} prints it as \texttt{\#RRGGBBAA} \\
\texttt{0x50} & 0 in user objects, 5 in system objects, 4 in the log and one data object \\
\texttt{0x58} & the length of DATA \\
\texttt{0x60} & where the table of relocations starts (0: none) \\
\bottomrule
\end{longtable}

\subsection{The parts, in their order}

\begin{enumerate}
\item \textbf{DATA}: the content as it is. A text, HTML, pixels, the PyTon text of a program.
      Empty in a compiled program.
\item \textbf{META}: texts \texttt{key:value}, each ended by a zero byte; an empty text ends the
      list. Keys seen: \texttt{type}, \texttt{name}, \texttt{label}, \texttt{kind}, \texttt{icon},
      \texttt{data}, \texttt{sockets}, \texttt{version}, \texttt{lang}, \texttt{source},
      \texttt{source\_hash}, \texttt{target}, \texttt{hash}, \texttt{author}, \texttt{screen},
      \texttt{buttons}, \texttt{rating}, \texttt{made}, \texttt{qr}.
\item \textbf{System sockets}: nine addresses and a zero. A socket the object does not have points
      at one byte \texttt{C3} (a bare \texttt{ret}) that follows the table. Which socket each of
      the nine places is, is not known yet.
\item \textbf{Named sockets}: one address for every name in META \texttt{sockets:}, in that order,
      then a zero. \cmd{run NAME [SOCKET]} calls one; the first is the default.
\item \textbf{Prockets}: sockets that fire by chance. An entry is an address, a percentage and one
      more number. Hello has one that fires half of the time.
\item \textbf{Code}: x86-64 machine code. A socket gets a pointer to the table of the system's
      functions and calls them through it (\texttt{call [r15+N]}).
\item \textbf{Relocations} (only in bigger programs): a count, then the places in the code the
      kernel corrects when it loads the object (the log: \texttt{relocated calc: 32 places}).
      A place with the top bit set is 32 bits wide, otherwise 64.
\item \textbf{The end}, 16 bytes: the size once more, then \texttt{ASMO} and four zero bytes.
      Gaps are filled with \texttt{0x90}.
\end{enumerate}

\subsection{Types of objects}

\begin{longtable}{@{}l p{11.5cm}@{}}
\toprule
\texttt{type:} & what it is \\ \midrule
\texttt{icon} & a program with an icon on the desktop (\texttt{kind:app}), or a launcher of
                something in the kernel (Photon, Puffer, Tuner, META: 560 bytes each) \\
\texttt{source} & the PyTon text of a program; \texttt{target} names the compiled object and
                \texttt{hash} ties the two together (\texttt{source\_hash} in the program) \\
\texttt{page} & HTML shown by NET \\
\texttt{style} & CSS: \texttt{theme} and the themes it is copied from (day, night, ocean) \\
\texttt{script} & a PyTon script a page runs \\
\texttt{image} & pixels, 32 bits each: icons 48$\times$48, pictures, the QR code \\
\texttt{font} & the letters of the system (Latin and Cyrillic) \\
\texttt{cabinet} & a drawer: room for 63 more objects (\texttt{archive}, \texttt{samples},
                \texttt{system}) \\
\texttt{kernel} & a kernel suit: \texttt{k\_tools}, \texttt{k\_usb12}, \texttt{k\_net},
                \texttt{k\_sound} \\
\texttt{desktop} & the desktop itself; its DATA is the wallpaper \\
\texttt{sentinel}, \texttt{machine}, \texttt{key}, \texttt{log}, \texttt{settings} &
                the list the Sentinel checks against, this machine's identity, the network key,
                the chained log, the settings \\
\texttt{demo}, \texttt{test}, \texttt{text}, \texttt{data}, \texttt{list} & small objects for
                showing and testing, and plain content \\
\bottomrule
\end{longtable}

% ------------------------------------------------------------------------------------------
\section{The disk}

The image of v2.9.59 is 46136 sectors of 512 bytes (22.5 MB).

\begin{longtable}{@{}l l p{9cm}@{}}
\toprule
sectors & bytes & what \\ \midrule
0 & \texttt{0x000000} & the boot sector \\
1 -- 480 & \texttt{0x000200} & the kernel: the core (61968 bytes), then the extension
                               (158032 bytes), then free room \\
from 490 & \texttt{0x03D400} & the cabinet: objects, each starting on a sector \\
\bottomrule
\end{longtable}

\noindent The first 0x3E bytes of the boot sector are the header of a FAT12 diskette (1.44 MB),
so that every BIOS takes the stick for a disk it can start. \textbf{Windows believes this header}:
it mounts a NinuxOS stick as FAT12 and writes \file{System Volume Information} into what it takes
for the directory and the first clusters, which are sectors 19, 44 and 45 of the kernel. A stick
that has been in Windows no longer starts properly. Until the header is changed: write sticks with
Puffus (it ejects them) and read them only with USB write protection on.

At the end of the boot sector the build (or \cmd{install}) leaves three numbers: at \texttt{0x1AC}
and \texttt{0x1B0} where the system starts on the disk (LBA, high and low half), at \texttt{0x1B4}
how many sectors the kernel takes (480).

The cabinet holds 165 objects: 105 that \cmd{ls} lists and 60 inside the drawers. \cmd{ls} prints
other addresses than the offsets in the image (they differ by \texttt{0xC4000}); how the system
keeps the list of objects, slots and gaps is not known yet.

% ------------------------------------------------------------------------------------------
\section{The start}

\begin{enumerate}
\item \textbf{The boot sector} (16-bit, at \texttt{0x7C00}) asks the BIOS whether it reads by LBA,
      takes the geometry for the other case, reads 480 sectors to \texttt{0x7E00} and jumps there.
      It leaves the drive number, the geometry and the start of the system at \texttt{0x500}.
\item \textbf{The loader} is the beginning of the core (16-bit, at \texttt{0x7E00}). It reads the
      memory map, picks the screen mode, finds the cabinet (``loading cabinet''), loads the GDT and
      switches to protected mode at \texttt{0x8360}.
\item \textbf{32-bit, shortly:} page tables at \texttt{0x18000} (2 MB pages), PAE, long mode, a jump
      to 64-bit code at \texttt{0x840B}. A second entry at \texttt{0x8423} is already 64-bit (the
      start from UEFI, v2.9.42).
\item \textbf{The core} runs where it was read, at \texttt{0x7E00}. It copies the extension to
      \texttt{0x50000} and checks it with SHA-256; if the hash is wrong the system starts with
      ``system operations off''.
\item Then, in the order of the log: PCI, ACPI, the boot disk, the font, the shield (GDT with user
      segments, the gates \texttt{0x80} and \texttt{0x81}, page tables that keep the kernel for the
      supervisor, a self-test in ring 3), the other processor cores, the screen, the memory, the
      interrupts (timer 100 Hz), the Sentinel, the suit \texttt{k\_tools}, the chained log.
\item The shell. After the desktop has run for ten seconds once, it starts by itself next time.
\end{enumerate}

% ------------------------------------------------------------------------------------------
\section{The kernel}

\subsection{The parts}

\begin{longtable}{@{}l l p{8.5cm}@{}}
\toprule
part & in memory & what it holds \\ \midrule
core & \texttt{0x7E00} & the loader, the start, memory, the gates, the shield, the cabinet, the
                         disk through the BIOS and ATA, the shell with its help text \\
extension & \texttt{0x50000} & begins with \texttt{NXEXT1} and its size; the drivers and the
                         bigger machinery: PCI, AHCI, NVMe, xHCI, virtio-gpu, the crypto, the
                         PyTon compiler, the Sentinel's work, the tuner, the SUB-GPU and Photon \\
suits & loaded when needed & objects of \texttt{type:kernel} in the cabinet: \texttt{k\_tools}
                         (the tools of the shell), \texttt{k\_usb12} (UHCI, OHCI),
                         \texttt{k\_net} (the network), \texttt{k\_sound} \\
\bottomrule
\end{longtable}

\noindent The core and the extension are linked as one: the extension calls functions of the core
directly, and the core calls the extension. The names of the old source files (Core, Memory, Gates,
Shield, Cart, ATA; ACPI, AHCI, Bench, Chain, Crypto, Files, Lock, NVMe, PCI, PhotonTables,
PyTonCompiler, Shot, SuitLoader, SysInfo, Tuner, USBStick, VirtioGPU, XHCI; the suits EHCI, OHCI,
UHCI, Sound) are known; which bytes belonged to which file is not.

\subsection{Protection}

\begin{itemize}
\item \textbf{The shield.} Objects run in ring 3 with their own memory and page tables (a new
      program gets its room at \texttt{0x1400000}); they reach the system only through the API
      (\texttt{int 0x80}) and the table of functions. A fault in an object is caught and reported
      (\texttt{[shield] object fault}); the system goes on. \cmd{shield on|off|test}.
\item \textbf{The Sentinel.} Every system object and the kernel (\texttt{@kernel}) are checked
      against the build's signed list (the object \texttt{sentinel}); \cmd{sentinel} prints
      \emph{genuine} or \emph{CHANGED}. A changed system object is refused until the administrator
      accepts it (\cmd{sentinel trust NAME}). Code the Sentinel has not verified runs only as a
      program in a window, never from the shell.
\item \textbf{Maintenance mode.} The cabinet can be changed only after \cmd{maint} and the
      administrator's password (\cmd{copy}, \cmd{rename}, \cmd{delete}, \cmd{build},
      \cmd{inject}, \cmd{install} \ldots).
\item \textbf{Identity.} A machine has an ID, a role, a generation, a parent and a key
      (\cmd{whoami}); a stick made by \cmd{puffer} is a child of the machine that made it.
      Objects can be locked with a passphrase, hidden, and rated by age (\cmd{level}, \cmd{rate}).
\item \textbf{The chained log} (\texttt{chainlog}) records what was done to the system.
\end{itemize}

\subsection{Programs side by side}

The \emph{multiplexer} runs the programs of the desktop in windows, each with its own memory and
time slices; the desktop is task 0. By default the programs take turns on the first processor core
(\cmd{cores all} lets them run on every core). Shared work --- drawing windows, the SUB-GPU,
Photon --- is spread over all cores.

\subsection{Devices}

Disks: the BIOS at the start, ATA, AHCI (SATA), NVMe; USB sticks on xHCI and EHCI. Input: PS/2,
USB keyboards and mice on xHCI, UHCI and OHCI. Screen: the framebuffer of the BIOS or of UEFI,
virtio-gpu. Sound: HD Audio and AC97. Network: e1000 and Realtek, with DHCP, and Meta Net on top
of it (\cmd{serve}, \cmd{get}, \cmd{share}, \cmd{send}, \cmd{subscribe}, \cmd{broadcast}).
Power: ACPI.

% ------------------------------------------------------------------------------------------
\section{PyTon}

A program is a text in the DATA of an object of \texttt{type:source}. It begins with a line of
META and defines sockets as functions:

\begin{verbatim}
META = {"type": "icon", "label": "Keys", "kind": "app", "icon": "icon_keys"}

def draw(x, y):         # the socket that draws the icon
    ...
def open(sw, sh):       # the socket that runs the program in its window
    ...
\end{verbatim}

\noindent The language has numbers, texts, lists, dicts, floats, functions, \texttt{global},
\texttt{if / elif / else}, \texttt{while}, \texttt{for}. The system gives functions such as
\texttt{fill}, \texttt{fill\_round}, \texttt{text}, \texttt{window}, \texttt{poll}, \texttt{key},
\texttt{mouse\_x}, \texttt{mouse\_y}, \texttt{mouse\_buttons}, \texttt{load\_text},
\texttt{save\_text}, \texttt{field}, \texttt{log}, \texttt{ob\_meta}. A program is an event loop:
\texttt{poll()} returns 1 for a key and 2 for the mouse.

\cmd{build} (maintenance mode) and F5 in Studio compile a source to a native object with
\texttt{lang:pyton}; \cmd{sources} lists each source, its target and whether the compiled program
is up to date. A new program comes in from a FAT32 stick with \cmd{files get NAME}.

Outside the system the same work was done by \file{pyton/py2asm.py}. That file is lost; its
ancestor (\file{py2asmx64v3.21.py}) is in the repository PypyOS.

% ------------------------------------------------------------------------------------------
\section{What the user sees}

\begin{itemize}
\item \textbf{The shell} (\texttt{ninux>}): about ninety commands; \cmd{help} lists them a page
      at a time. The name of an object alone runs it.
\item \textbf{The desktop} (\cmd{desktop}): icons of the objects, a taskbar, the start menu, a menu
      on the right button; a drawer opens as a desktop of its own. Esc closes the program on top;
      Esc on the bare desktop goes back to the shell.
\item \textbf{The programs} of v2.9.59: Ask, ASMONCRAFT, Bags, Calc, Chronos, Console, Copycat,
      Device, Greet, Interfacer, Keys, Mail, Nest, NET, Notepad, Paint, Phone, Projector, PyNamp,
      Search, Studio, Widgets; full screen: Photon, Tuner, the META launcher; and Puffer.
\end{itemize}

% ------------------------------------------------------------------------------------------
\section{The files of the project}

\begin{longtable}{@{}p{4.6cm} p{10.4cm}@{}}
\toprule
path & what it is \\ \midrule
\file{disk.img} & the system, ready to start: what the build writes and Puffus puts on a stick \\
\file{Puffus.exe} & Puffus for Windows (made from \file{tools/puffus}) \\
\file{CHANGELOG.md}, \file{TODO.md} & what each version brought; what is left to do \\
\midrule
\file{boot/Boot.kasm} & the boot sector, NASM source with names and comments \\
\file{kernel/Kernel.kasm} & the kernel (core and extension) as NASM source, made by
      \file{tools/kdis}. Labels are addresses: \texttt{fn\_} a function, \texttt{L\_} a place
      jumped to, \texttt{d\_} data. About half is recognised as code; the rest is \texttt{db}.
      \textbf{Lines must not be added or removed yet}: addresses in data are still numbers \\
\file{apps/*.asmo} & the PyTon text of the programs (plain text, despite the name) \\
\file{pages/*.asmo} & the pages (HTML) and the themes (CSS) \\
\file{scripts/*.asmo} & the PyTon scripts of the pages \\
\file{texts/} & objects of plain text without a name, and the log of the system \\
\midrule
\file{build/objects/*.asmo} & every object as it lies in the cabinet (header, DATA, META, sockets,
      code); \file{X.Source.asmo} is the source object of program X \\
\file{build/objects/Kernel.kasm} & the compiled kernel, 220000 bytes; \file{Boot.bin} the compiled
      boot sector \\
\file{build/kernel/*.kasm} & the compiled kernel suits: Tools, USB12, Network (\texttt{k\_net}),
      Sound \\
\file{build/base/disk\_v2.9.59.img} & the system exactly as it came off the stick; the build
      starts from it \\
\midrule
\file{tools/build\_disk.py} & the build: assembles boot and kernel, puts them and the objects
      into the base image, writes \file{disk.img} \\
\file{tools/puffus/} & \file{Puffus.cs}, its manifest and \file{make\_puffus.cmd} \\
\file{tools/kdis/kdis.py} & makes \file{kernel/Kernel.kasm} from the compiled kernel and proves
      the result assembles to the same bytes; \file{trace\_addrs.txt} are addresses seen running
      in QEMU; optional \file{names.txt} and \file{seeds.txt} give names and more code starts \\
\file{tools/qemu/qrun.py} & starts the system in QEMU and drives it from a script \\
\file{tools/dump\_usb.ps1} & reads a stick into an image, sector by sector (Administrator) \\
\file{test/qemu/q\_*.txt} & the test runs: the shell, the desktop, the programs \\
\midrule
\file{docs/NinuxOS\_Architecture.tex} & this document \\
\file{docs/cabinet\_objects.csv} & every object: its offset in the image, size, type, name, its
      file in \file{build}, its readable source \\
\file{docs/chat/} & the conversation in which the project was brought back \\
\bottomrule
\end{longtable}

\subsection{From source to a stick}

\begin{enumerate}
\item \cmd{python tools/build\_disk.py} runs NASM on \file{boot/Boot.kasm} and
      \file{kernel/Kernel.kasm}, checks their sizes, puts them and every object of \file{build}
      into a copy of the base image and writes \file{disk.img}. It says what differs from v2.9.59.
\item \file{Puffus.exe} (as Administrator) runs the build, lists the USB disks, locks and
      dismounts the chosen stick, writes \file{disk.img} from its first sector, reads it back,
      compares, ejects.
\item To try without a stick: \cmd{python tools/qemu/qrun.py SCRIPT.txt --show}; the image is
      opened so that nothing is written into it.
\end{enumerate}

\noindent Needed on the PC: Python 3, NASM, QEMU; for Puffus only the .NET Framework that
Windows has.

% ------------------------------------------------------------------------------------------
\section{What is not known yet}\label{sec:open}

\begin{itemize}
\item How the identity of an object (header \texttt{0x4C}) is computed.
\item Where the expected SHA-256 of the extension lies, and the format and the key of the Sentinel's
      list. Until then a kernel or a system object with other bytes is refused, so the system in
      v2.9.60 is still the v2.9.59 kernel.
\item The table of the cabinet: slots, drawers, gaps; and 580 KB between the objects of the image
      that belong to none of them.
\item The nine system sockets and the third number of a procket; the fields \texttt{0x48} and
      \texttt{0x50} of the header.
\item The numbers of the functions in the API table.
\item Which part of the kernel was which source file.
\end{itemize}

\end{document}
